\section{案例阶梯：细胞、类器官、结构与工具反馈}

课程接着展示多种合作案例。它们不应被揉成一句“AI 发现了很多药物和蛋白质”，因为每个案例的研究对象、实验系统、发表状态和可外推范围不同。本章按证据等级分别阅读 AML、肝纤维化、植物免疫组装、OCT4 设计和早期 rejuvenation 结果。

\subsection{先建立案例地图}

总览图把工作分成癌症药物再利用、逆转肝损伤，以及发现和工程化蛋白质。共同流程是 AI 提出或组织候选，人类实验团队选择验证；不同之处在于验证平台从细胞系、类器官到结构生物学不等，因而结论强度不能横向直接比较。

\lecturefigure{slide-009.jpg}{人类--AI 协作正在推进的案例地图}{Stanford CS25 V6 Lecture 07 官方录像 00:42:00}

\begin{knowledgebox}{读图：三个方向共享流程，不共享结论强度}
左侧 AML 主要看药物对癌细胞系 viability 的影响；中间肝纤维化使用人源 hepatic organoid 观察抗纤维化与再生信号；右侧蛋白质方向把 AlphaFold 结构筛选接到纯化和显微实验。它们都可证明“候选值得继续研究”，却分别距离临床治疗、器官级疗效和生物功能确认还有不同距离。
\end{knowledgebox}

\begin{importantbox}{案例状态总表}
\begin{tabularx}{\textwidth}{L{0.20\textwidth}X X}
\toprule
案例 & 课堂日前证据 & 讲义允许的结论 \\
\midrule
AML / KIRA6 & arXiv v1 中的体外细胞系结果 & 有选择性 cytotoxicity 候选，需进一步临床前验证 \\
Liver fibrosis & bioRxiv v1，人源 hepatic organoid & 两类 epigenomic modifier 有抗纤维化和再生信号 \\
Plant NRC7 & bioRxiv v1，AlphaFold 筛选、纯化、负染电镜 & SNI 找到非典型 11-mer assembly \\
OCT4 & arXiv v1 appendix 的结构 plausibility demonstration & 结构未明显破坏，功能仍未证实 \\
Rejuvenation & 课堂 slide，明确 unpublished & 仅记录初步课堂结果，不作定论 \\
\bottomrule
\end{tabularx}
\end{importantbox}

\subsection{AML：KIRA6 的信号与两个失败候选}

急性髓系白血病（acute myeloid leukemia, AML）药物再利用任务要求系统从已有药物中寻找新适应证。课堂重点是 KIRA6：一种 IRE1$\alpha$ inhibitor，在若干 AML cell line 上出现剂量依赖的 viability 下降。图也保留了关键负面信息：AI 提议的另外两个“新”药没有工作。

\lecturefigure{slide-011.jpg}{KIRA6 在 AML 细胞系上的剂量--反应结果}{Stanford CS25 V6 Lecture 07 官方录像 00:45:28；arXiv:2502.18864v1}

\begin{knowledgebox}{读图：先看横轴浓度，再看不同细胞系的下降位置}
横轴是对数浓度，纵轴是相对 viability。曲线越早下降，表示较低浓度即可抑制细胞活性；但不同细胞系的敏感性并不相同。真正有价值的是把 AML cell line 与非恶性对照放在同一设计中寻找 therapeutic window，而不是只展示一条下降曲线。图中的 novelty review 也承认已有相关靶点证据，novelty 在于具体药物和适应证组合。
\end{knowledgebox}

常用 Hill dose--response 形式为
$$
V(c)=V_{\min}+\frac{V_{\max}-V_{\min}}
{1+(c/IC_{50})^{h}}.
$$
其中：
\begin{itemize}
  \item $c$ 是药物浓度；
  \item $V(c)$ 是该浓度下测得的细胞活性；
  \item $IC_{50}$ 是活性下降到中间水平所需浓度；
  \item $h$ 是 Hill slope，描述曲线陡峭程度；
  \item 体外 $IC_{50}$ 不能直接替代人体剂量、安全性或疗效。
\end{itemize}

\teachervoice{00:45:20--00:47:12，老师主动指出两种其他新药建议没有工作。负结果不是演讲瑕疵，而是估计系统命中率、发现 selection bias 和设计下一轮 ranking 的必要数据。}

\begin{warningbox}{Cell killing 不是临床治疗}
细胞系缺少人体中的药代动力学、免疫环境、组织分布、毒性和耐药演化。KIRA6 的结果支持继续做机制、动物与安全研究，不支持把它描述成 AML 已获证实的新疗法。
\end{warningbox}

\subsection{肝纤维化：人源类器官中的 Vorinostat}

肝纤维化案例把验证平台升级为 multi-lineage human hepatic organoid。AI co-scientist 推荐 epigenomic modifier，实验团队串行评估抗纤维化活性与毒性。Vorinostat 是已获批的抗癌 HDAC inhibitor，在该模型中降低 TGF$\beta$ 诱导的染色质结构变化，并伴随抗纤维化与肝实质细胞再生信号。

\lecturefigure{slide-012.jpg}{Vorinostat 在人源 hepatic organoid 中的抗纤维化结果}{Stanford CS25 V6 Lecture 07 官方录像 00:47:48；bioRxiv:2025.04.29.651320v1}

\begin{knowledgebox}{读图：虚线是基准，不是治愈线}
纵轴是 fibroblast activity 相对 untreated 的 fold change，绿色虚线约为 1。Fibrosis inducer 把活性推高；若候选药把分布拉回或低于基准，说明在该 organoid 条件下减少促纤维化活动。箱线图展示样本分布而非单个最佳点。Vorinostat 的价值在于一个已知药物获得新机制和新适应证候选，但图不包含人体结局。
\end{knowledgebox}

若用 untreated 基准归一化，抗纤维化效应可写为
$$
E_{\text{anti-fibrotic}}
=1-\frac{A_{\text{drug}}-A_{\text{baseline}}}
{A_{\text{fibrosis}}-A_{\text{baseline}}}.
$$
其中：
\begin{itemize}
  \item $A_{\text{drug}}$ 是候选药处理后的纤维化活动指标；
  \item $A_{\text{fibrosis}}$ 是诱导纤维化但未给候选药的指标；
  \item $A_{\text{baseline}}$ 是未诱导基准；
  \item 该归一化有助于比较实验内效果，不能消除模型与人体之间的外推误差。
\end{itemize}

\begin{knowledgebox}{背景概念：organoid}
类器官是由细胞在三维环境中自组织形成、复现部分组织结构和功能的实验模型。它比单层细胞系更接近组织级相互作用，却通常仍缺少完整血流、免疫、神经和全身代谢。因此 organoid 证据强于纯文本推理，仍早于动物、临床前和人体试验。
\end{knowledgebox}

\teachervoice{00:47:20--00:48:46，老师把实验平台视为对 AI 假设的 biological reality check。课程没有把 organoid 改善包装成病人治疗成功。}

\subsection{植物结构：SNI 怎样把异常变成候选}

AlphaFold 生成大量结构预测后，研究者面临另一个问题：哪些“不像已知结构”的结果是真正新颖，哪些只是预测错误？团队与 AI co-scientist 设计 Structural Novelty Index（SNI），对植物免疫蛋白组装进行结构筛选，并把异常候选交给纯化和负染 electron microscopy。

\lecturefigure{slide-013.jpg}{Structural Novelty Index 找到 NRC7 的非典型 11-mer assembly}{Stanford CS25 V6 Lecture 07 官方录像 00:50:32；bioRxiv:2026.05.03.722499v1}

\begin{knowledgebox}{读图：AI 负责缩小搜索，人类实验决定异常是否真实}
上方热图比较科学家与 AI co-scientist 对候选的结构新颖性判断；下方流程从 SNI 筛选进入表达、纯化与 electron microscopy，最终观察到 11 个亚基组成的环状 NRC7 assembly。教学重点不是“AI 看图比人强”，而是它提出额外判据，帮助从数百个 AlphaFold 结构中优先选择值得昂贵实验检查的异常。
\end{knowledgebox}

\begin{importantbox}{Anomaly detection 的科研价值}
当生成模型产出海量候选时，最稀缺的不是候选，而是实验注意力。SNI 一类指标把“与 canonical architecture 的偏离”转成可排序信号；随后必须用独立实验区分生物学 novelty 与模型 artifact。Agent 最有价值的角色常是缩小验证空间，而不是宣布异常已经成立。
\end{importantbox}

\subsection{OCT4：把 AlphaFold 放进反馈环}

OCT4 是 Yamanaka factors 之一。AI co-scientist 提议修改其 POU domain 附近的序列，再调用 AlphaFold 检查结构 plausibility，并把结构结果反馈到下一轮设计。arXiv v1 appendix 明确称其为 illustrative demonstration：若要证明改造提高 DNA binding 或保留 SOX2 interaction，还需要 binding affinity、off-target 与湿实验验证。

\lecturefigure{slide-014.jpg}{AI co-scientist 与 AlphaFold 组成 OCT4 设计反馈环}{Stanford CS25 V6 Lecture 07 官方录像 00:51:16；arXiv:2502.18864v1 Appendix A.6}

\begin{knowledgebox}{读图：pLDDT、pTM、ipTM 只回答结构置信问题}
图比较原始 OCT4 与修改序列的预测结构和置信指标。相近 pLDDT 表明局部结构置信度未明显崩溃；pTM 与 ipTM 分别提供整体折叠和界面可信度线索。它们不能直接证明转录功能增强、细胞命运改变或没有脱靶。正确结论是“结构上值得继续验证”，不是“优化蛋白已经成功”。
\end{knowledgebox}

\begin{lstlisting}[language=Python,caption={专门模型进入假设设计反馈环的概念流程}]
sequence = co_scientist.propose_sequence(goal, constraints)
for iteration in range(max_rounds):
    structure = alphafold.predict(sequence)
    checks = assess(structure, metrics=["pLDDT", "pTM", "ipTM"])
    if checks.structurally_plausible:
        sequence = co_scientist.refine(sequence, checks)
    else:
        sequence = co_scientist.repair(sequence, checks.failures)
return send_to_binding_and_wet_lab_validation(sequence)
\end{lstlisting}

\begin{warningbox}{Tool use 不是 tool truth}
调用 AlphaFold 能减少明显不合理序列，却不能把预测模型变成实验仪器。若 Agent 与工具共享偏差，反馈环可能稳定优化错误代理指标。必须在 loop 外加入异构模型、数据库检查、实验和停止条件。
\end{warningbox}

\subsection{Early rejuvenation：课堂结果必须保留红线}

课程展示了老年供体细胞中 senescence marker 大幅下降的早期结果，并在标题中直接写明 `Unpublished and needs peer review`。这类图很容易制造“逆转衰老已实现”的印象，因此讲义必须同时记录实验对象、marker、对照以及尚未公开方法和复现的事实。

\lecturefigure{slide-015.jpg}{新型 rejuvenation factor 的早期课堂结果}{Stanford CS25 V6 Lecture 07 官方录像 00:52:00；slide 明确标注 unpublished and needs peer review}

\begin{knowledgebox}{读图：柱高描述 marker，不等于年龄倒流}
纵轴是被判定为 senescent 的细胞比例，老年供体 untreated 组较高，KLOTHO 和两个 AI-discovered factor 组明显降低。图支持“在该实验条件下 marker 下降”的初步观察；它没有说明细胞所有功能恢复年轻状态，也没有给出样本量、完整统计、脱靶效应和独立重复。
\end{knowledgebox}

\begin{warningbox}{Unpublished means what}
未发表不等于错误，也不等于已证实。它意味着读者无法完整核查方法、样本、排除标准、统计分析和复现。讲义保留这个案例，是为了学习证据分级与研究方向，不将其写成抗衰老疗法结论。
\end{warningbox}

\teachervoice{00:51:40--00:53:32，老师在 slide 上主动标出 needs peer review。教师语气是“early results”，不是正式发布；任何更强措辞都会越过课堂证据。}

\subsection{本章小结}

五个案例构成从候选生成到不同实验层级的阶梯。AML 提供细胞系 dose--response 并暴露失败候选；肝纤维化进入人源类器官；植物 NRC7 用结构筛选加纯化和电镜确认 11-mer；OCT4 只是结构 plausibility feedback；rejuvenation 仍是未发表课堂结果。Agent 可以提高候选发现和实验优先级，但每种实验模型都有明确外推边界。

